\documentclass[conference]{IEEEtran}
\IEEEoverridecommandlockouts
% Force Arabic numbering for sections and subsections as requested
\renewcommand\thesection{\arabic{section}}
\renewcommand\thesubsection{\thesection.\arabic{subsection}}
\renewcommand\thesubsubsection{\thesubsection.\arabic{subsubsection}}

\usepackage{cite}
\usepackage{amsmath,amssymb,amsfonts}
\usepackage{algorithmic}
\usepackage{graphicx}
\usepackage{textcomp}
\usepackage{xcolor}
\usepackage{hyperref}
\usepackage{orcidlink}
\hypersetup{hidelinks}
\def\BibTeX{{\rm B\kern-.05em{\sc i\kern-.025em b}\kern-.08em
    T\kern-.1667em\lower.7ex\hbox{E}\kern-.125emX}}

\begin{document}

\title{OmniGen: Evaluation Against SOTA Image Generation Models}

\author{\IEEEauthorblockN{Huy Ngoc Chau \orcidlink{0009-0004-3554-059X}}
\IEEEauthorblockA{\textit{University of Science} \\
\textit{Viet Nam National University}\\
Ho Chi Minh City, Vietnam \\
25C1100828@student.hcmus.edu.vn \\}
\and
\IEEEauthorblockN{Minh Thang Hoang \orcidlink{0009-0004-3554-059X}}
\IEEEauthorblockA{\textit{University of Science} \\
\textit{Viet Nam National University}\\
Ho Chi Minh City, Vietnam \\
25C1102003@student.hcmus.edu.vn \\}
}

\maketitle

\begin{abstract}
OmniGen is a unified image generation model designed to handle multiple image generation tasks within a single framework. It is among the first approaches to jointly support diverse tasks, while achieving performance comparable to task-specific state-of-the-art (SOTA) models.

Motivated by these claims, we conduct a systematic evaluation of OmniGen against recent SOTA models across three representative tasks: text-to-image generation, instruction-based image editing, and style transfer. Our evaluation considers both task-specific performance metrics and computational efficiency.

Experimental results, consistent with those reported in the original work, indicate that OmniGen achieves competitive performance in general image generation tasks. However, it shows limitations in tasks requiring precise control and fine-grained manipulation, where specialized models still outperform the unified approach.

These findings highlight the trade-offs between generality and task-specific performance, suggesting that while unified models offer flexibility and simplicity, specialized architectures remain essential for achieving optimal results in individual tasks.
\end{abstract}

\begin{IEEEkeywords}
OmniGen, Reproducibility, Unified Image Generation, Diffusion Transformer, Multimodal Generation
\end{IEEEkeywords}

\section{Introduction}
\label{sec:introduction}

The development of large-scale generative models has transformed many domains of artificial intelligence. In natural language processing, large language models handle diverse tasks in a unified framework. In visual generation, however, practical workflows are still fragmented across multiple task-specific models.

Existing systems are typically specialized: text-to-image models such as Stable Diffusion are strong for prompt-based synthesis \cite{rombach2022high}, while methods such as ControlNet provide conditional control through dedicated branches \cite{controlnet}. This fragmentation increases engineering overhead because complex workflows require chaining different models, each with separate preprocessing and inference settings.

OmniGen \cite{omnigen2024} proposes a unified alternative: a single diffusion-transformer framework that accepts interleaved text and image inputs for multiple tasks, including text-to-image generation, instruction-based editing, and subject-driven generation. Motivated by this claim, we conduct a systematic reproduction and evaluation study on representative benchmarks and compare against state-of-the-art baselines.

Beyond benchmark-level analysis, we also study practical deployment by integrating OmniGen into an AI-powered comic generation pipeline. This application perspective helps identify real-world limitations that may not be fully visible in standard benchmark scores.

Our contributions are:
\begin{enumerate}
    \item A reproducible evaluation protocol for OmniGen on GenEval, EMU-Edit, and DreamBench.
    \item A concise technical summary of architecture and training elements that affect inference behavior.
    \item A practical assessment through comic generation, including observed limitations in text rendering, consistency, multilingual prompting, and multi-panel detail preservation.
\end{enumerate}

The remainder of this paper is organized as follows. Section \ref{sec:related_work} reviews prior work. Section \ref{sec:method} describes OmniGen and our reproduction protocol. Section \ref{sec:experiments} reports results and discussion. Section \ref{sec:conclusion} concludes with key insights and future directions.

\section{Related Work}
\label{sec:related_work}

\subsection{Text-to-Image Generation}
Diffusion-based text-to-image models have achieved major progress in visual quality and compositional alignment. Stable Diffusion and its successors demonstrated efficient latent-space denoising \cite{rombach2022high}, while large diffusion transformers such as SD3 improved quality at scale \cite{sd3}. These models are highly capable, but they are still primarily optimized for prompt-to-image synthesis.

\subsection{Controllable and Multi-task Generation}
To support finer control, prior work introduced task-specific modules. ControlNet augments base diffusion models with condition branches for edge, depth, and segmentation control \cite{controlnet}. Editing and identity-preserving methods similarly rely on dedicated components, often producing strong task-specific performance but requiring multi-stage pipelines and additional model orchestration.

\subsection{Unified Image Generation Models}
Recent research explores unifying multiple generation tasks under shared backbones. OmniGen is a representative approach that removes separate condition encoders and instead uses a single multimodal transformer with interleaved text-image instructions \cite{omnigen2024}. This design motivates direct comparison between unified and specialized paradigms.

\subsection{Evaluation Benchmarks}
For fair comparison, the community commonly uses GenEval for compositional text-to-image alignment \cite{geneval}, EMU-Edit for instruction-based editing \cite{emuedit}, and DreamBench for subject-driven generation \cite{dreambench}. We adopt these benchmarks to maintain comparability with prior studies.

\section{Method}
\label{sec:method}

\subsection{OmniGen Architecture Summary}
OmniGen consists of two main modules: a frozen VAE and a transformer backbone initialized from Phi-3 \cite{phi3,omnigen2024}. The VAE maps images to latent representations and decodes generated latents back to pixel space. Text inputs are tokenized by the language tokenizer, while image latents are converted into patch tokens and inserted into a unified sequence with image boundary tokens.

The model uses mixed attention: causal attention across the full multimodal sequence and bidirectional attention within each image token block. This structure preserves autoregressive instruction flow while enabling stronger local visual interaction among patches from the same image.

At inference time, OmniGen samples Gaussian noise in latent space and performs iterative velocity prediction with rectified flow, then decodes the final latent through the VAE decoder.

\subsection{Training Objective}
OmniGen follows rectified flow training \cite{omnigen2024,sd3}. At diffusion time $t$, the noisy latent is
\begin{equation}
\mathbf{x}_t = t\mathbf{x} + (1-t)\boldsymbol{\epsilon}, \quad \boldsymbol{\epsilon} \sim \mathcal{N}(0,1)
\end{equation}
and the velocity prediction objective is
\begin{equation}
\mathcal{L} = \mathbb{E}\left[\left\|(\mathbf{x}-\boldsymbol{\epsilon}) - v_\theta(\mathbf{x}_t,t,c)\right\|_2^2\right].
\end{equation}

For editing tasks, OmniGen applies a region-weighted loss to reduce trivial identity-copy behavior from source to target image. Regions with larger source-target differences receive higher optimization emphasis, improving instruction-following edits.

\subsection{Reproduction Protocol}
We reproduced OmniGen using the official public implementation and released checkpoints \cite{omnigen2024}. No fine-tuning was performed. Following the original setup, we used 50 inference steps, text guidance scale $2.5$, image guidance scale $1.6$, and image editing guidance scale $2.0$.

Evaluation follows official benchmark protocols for GenEval \cite{geneval}, EMU-Edit \cite{emuedit}, and DreamBench \cite{dreambench}. To ensure consistency across tasks, all prompts were converted into OmniGen's interleaved multimodal instruction format.

For downstream application analysis, we integrated OmniGen into a comic generation pipeline that includes translation, storyboard generation, panel synthesis, and text overlay post-processing.

\section{Experiments}
\label{sec:experiments}

\subsection{Experimental Setup}
We evaluate OmniGen on representative tasks aligned with its unified claims: text-to-image generation on GenEval \cite{geneval}, instruction-based editing on EMU-Edit \cite{emuedit}, and subject-driven generation on DreamBench \cite{dreambench}. We additionally report controllable generation quality under segmentation, Canny edge, HED edge, and depth conditions, following common ControlNet-style protocols \cite{controlnet}.

All experiments were run with a memory-efficient inference configuration on a single NVIDIA A100 80GB GPU. Output resolution was fixed at $1024 \times 1024$ for benchmark comparisons.

\subsection{Evaluation Metrics}
We report the GenEval overall compositional score for text-to-image alignment. For EMU-Edit and DreamBench, we report CLIP-T and CLIP-I using CLIP embeddings \cite{radford2021learning}. For controllable generation, we report mIoU (segmentation), F1 score (Canny), SSIM (HED), and RMSE (depth).

\subsection{Results}
Table \ref{tab:generation_results} reports benchmark outcomes. OmniGen reaches a GenEval overall score of $0.70$, and remains competitive on EMU-Edit and DreamBench, indicating that a unified model can approach specialized systems across heterogeneous tasks.

\begin{table}[htbp]
\caption{Key Benchmark Results for OmniGen}
\begin{center}
\begin{tabular}{|c|c|c|}
\hline
\textbf{Benchmark} & \textbf{Metric} & \textbf{OmniGen} \\
\hline
GenEval & Overall & 0.70 \\
\hline
EMU-Edit & CLIP-T / CLIP-I & 0.231 / 0.829 \\
\hline
DreamBench & CLIP-T / CLIP-I & 0.315 / 0.801 \\
\hline
Control Conditions & Seg. mIoU & 40.06 \\
\hline
\end{tabular}
\label{tab:generation_results}
\end{center}
\end{table}

For controllable generation, OmniGen performs strongly across multiple condition types and achieves best-in-group values on several metrics reported in the original benchmark comparison. Qualitatively, the unified prompting format supports end-to-end editing, subject transfer from references, and condition-following generation without switching models.

\subsection{Discussion}
Our reproduction is broadly consistent with the original OmniGen report \cite{omnigen2024}. Results support the claim that unified generation can remain competitive with specialized baselines while simplifying deployment.

However, practical application reveals limitations that are important for reproducibility and real-world use:
\begin{itemize}
    \item \textbf{Weak in-image text rendering}: generated typographic content is often distorted or unreadable.
    \item \textbf{Limited cross-panel consistency}: recurring characters may drift in appearance across comic panels.
    \item \textbf{English prompt dependency}: performance degrades with non-English prompts in our tests.
    \item \textbf{Detail loss in multi-panel layouts}: generating many panels jointly reduces per-panel visual fidelity.
\end{itemize}

These findings suggest a clear trade-off: unified models improve usability and flexibility, while specialized modules still provide stronger control for high-precision production workflows.

\section{Conclusion}
\label{sec:conclusion}
This paper presents a systematic reproduction and evaluation of OmniGen, a unified image generation framework. Across GenEval, EMU-Edit, and DreamBench, OmniGen demonstrates competitive results relative to task-specific baselines while maintaining a single-model interface.

Our downstream comic-generation study complements benchmark evaluation by exposing practical bottlenecks in text rendering, visual consistency, multilingual prompting, and multi-panel detail preservation. These observations highlight the core trade-off between generality and specialized control.

Future work includes domain-specific fine-tuning for narrative illustration, stronger identity-preserving mechanisms, and training strategies that improve typography and high-detail multi-panel generation.

% Keep bibliography non-empty even when all in-text citations are commented out.
\nocite{omnigen2024}

\bibliographystyle{IEEEtran}
\bibliography{references}

\end{document}
